\documentclass[12pt]{article}

\usepackage[a4paper, total={6in, 8in}]{geometry}
\usepackage{textcomp}

\begin{document}
\setlength{\parindent}{0pt}

\def \t {\textrightarrow}
\def \v {\vspace{1em}}
\def \bi {\begin{itemize}}
\def \ei {\end{itemize}}
\def \s[#1] {\section*{#1}}
\def \ss[#1] {\subsection*{#1}}
\def \sss[#1] {\subsubsection*{#1}}

\s[ONU]
Si vogliono trovare degli accordi che generarano strutture permanenti che rappresentino gli interessi di tutte le nazioni \t questo era anche l ideale della societa delle nazioni, che non funziona per contesto diverso e le sanzioni erano sottodimensionate rispetto agli eventi \\
Questo è il secondo tentativo \t l onu nasce dall'esigenza di garantire la pace, e un luogo dove le nazioni potessere parlarsi in modo diplomatico \\
L'onu è stata la prima grande org. internazionale \t dopo onu queste organizzazioni si sono moltiplicate, perche onu crea le agenzie come UNESCO, UNICEF, etc. \\
Sono strutture che fanno riferimento all onu, ma si comportano di argomenti specifici \t facendo riferimento all onu ma sono autonome \\
Sono organizzazioni a carattere universale \t includono potenzialmente tutte le nazioni del mondo \\
Poi ci sono serie di organizzazioni internazionali che non sono mondiali, a carattere regionale \t si riferiscono a stati di una parte specifica del pianeta 

\v

L'onu nasce nel 45 con la firma di 51 paesi sotto lo statuto delle nazioni unite \t ha sede a new york, e tra anni 50 e 60 sono entrati molti altri paesi (che arrivano dal processo di decolonizzazione) \\
L'altra crescita si ha negli ani 90 \t con dissoluzion urss paesi dell est diventano paesi sovrani \\
Attualmente onu comprende tutti i paesi indipendenti del mondo tranne:
\bi
\item vaticano (che pero ha a new york osservatore permanente)
\item taiwan \t fino al 71 sedeva nel consiglio di sicurezza, ma con mao viene allontana al posto della cina
\item la palestina è presente all onu con osserbatore (come vaticano) dal 2012
\ei

\v

Lo statuto definisce struttura e i suoi compiti:
\bi
\item mantenimento della pace nel mondo = risoluzione diplomatica delle controversie, e quindi:
  \bi
    \item difendere uguaglianza dei paesi memebri
    \item paesi membri sono indipendenti e devono essere difesi da aggressioni
    \item 
  \ei
\item favorire cooperazione tra popoli, ovvero mantenere il rispetto dei diritti e delle libertà degli individui \t non imporre regole 
\ei

\v

ONU è composta da:

\ss[Assemblea generale]
Si riunisce una volta all anno in sessione ordinaria, ma puo riunirsi anche con sessione straordinaria per questioni urgenti \\
Nell'assemblea ogni stato ha un voto, indipendemente dalla grandezza richezza e importanza di uno stato \\
L'assemblea puo prendere due tipi di decisioni:
\bi
  \item deliberare questioni importanti, ovvero quelle che per essere risolte richiedono maggioranza qualificata (2 terzi assemblea)
    \bi
      \item ammissione all onu di altri membri, o esplusione
      \item raccomandazioni sul mantenimento della pace a uno stato che le sta mettendo in pericolo \t puo richiamare uno stato che sta la mette in pericolo, pero non sono vincolanti queste racc. \t il non seguirle pero mette il paese in una situazione internazionale diversa
      \item l'onu richiama uno stato richiamandolo a dei principi fondamentali \t e alla fine gli stati hanno assunto comportamento delle raccomandazioni \t es. 1948 dichiarazione dei diritti fondamentali: onu espresso ha posizioni di principio, che poi ha influenzato azione dei paesi
    \ei
  \item deliberare altre questioni, ovvero adottate anche a maggioranza sempli (51 per cento)
\ei

\ss[Consiglio di sicurezza]
15 membri, di cui 5 permanenti e 10 non \t i non permanenti vengono eletti ogni 2 anni dall'assemblea generale \\
I 5 membri sono:
\bi
  \item francia
  \item uk
  \item usa
  \item russia
  \item cina
\ei
E questi paesi hanno il diritto di veto \t che ha reso difficile l'operato dell onu \\
Nel consiglio ogni membro ha un voto, per prendere decisioni meno importanti è necessaria maggioranza di 9 voti \t mentre per le decisioni importanti serve unanimità e consenso di tutti i membri permantenti \t che sono diversi ideologicamente \\
Quindi le decisioni possono essere bloccate da un unico voto, quindi possono bloccare il lavoro dell onu

\v

Il consiglio di sicurezza adotta le risoluzioni dell'onu \t che devono essere rispettate, e se non lo sono onu puo adottare misure economiche \\
Risoluzioni sono decisioni \t se stati non le rispettano, consiglio di sicurezza puo fare in modo che i paesi membri non commerciano con quel paese, sanzioni o addirittura intervento militare \\
Onu non ha esercito proprio \t e conta sulle forze armate messe dagli stati membri \t le forze armate dell onu dovrebbero in realta essere forze di pace \t i caschi blu devono andare nel paese per portare la pace

\ss[Segretariato generale]
Il segretario generale è il rappresentante dell onu nel mondo \t è la figura + visibile dell onu e la + autorevole \\
Viene nominato dall'assemblea generale su propsota del consiglio di sicurezza, carica di 5 anni (rinnovabili) \t è il funzionario piu alto in grado dell onu \\
Ha il compito di fare applicare concretamente le risoluzini del onu \t fa grande lavoro diplomatico e di mediazione \t le delibere dell onu non si possono imporre con la forza \\
Il segretario generale è anche a capo dell apparato burocratico dell onu, e di tutte le org. che lavorano per onu nei paesi membri

\ss[Corte internazionale di giustizia]
Ha sede all Aia in paesi bassi \t deve risolvere giuridicamente conflitti \\
Fornisce risposte a questioni giuridiche complesse che onu le sottopone \\
È composta da 15 giudici, che restano in carica 9 anni \t e nominati da assemblea generale e consiglio sicurezza \\
I giudici vengono scelti con il criterio di rappresentare tutti i sistemi giuridici nel mondo \t è un tribunale equilibrato, senza una scelta di sistema giuridico specifico \t ma equilibirio \\
Le sentenze sono vincolanti, e lo stato destinatario della sentenza se non si adegua puo essere "perseguito" \t gli puo essere imposto pena l'espulsione dall onu \\
Milosevic è stato giudicato dal tribunale dell Aia

\ss[consiglio economico e sociale]
Rappresentanti di 54 stati membri, eletti da assemblea generale \\
Aiuta lo sviluppo delle zone meno avanzate del pianeta \t e risolve emergenze alimentari e sanitarie

\s[Unione Europea]
UE ha storia antica e ha radici molte italiane \t gia De Gasperi ne parla \\
UE nasce nel 92 con trattato di Mastricht (in Olanda) \t gli stati che facevano parte della CEE firmano il trattato di M. \\
L UE è un organizzazione internazionale regionale \t è nata e funziona su scala continentale, a differenza dell onu \\
Viene fondata da 12 membri originari, che diventano 15 nel 95, e poi si arriva a 28 paesi \t e come all onu, nuovi stati che nascono durante dissoluzione dell urss chiedono di entrare in UE \\
Polonia, slovenia, etc. entrano nel 2004 \\
L UE è cresciuta e comprende 500 milioni di cittadini nei suoi confini \t e negli anni 90 e 2000, l UE aveva un PIL che era al secondo posto dopo la cina (ora non piu)

\ss[Storia]
UE è nata con una matrice economica \t CEE nasce da un mercato, e sul piano economico UE ha fatto progressi, e l unione dei sistemi dei paesi è stata raggiunta \t e paesi membri si impongono regole sulla regolazione della loro economia \\
Dal punot di vista economico funziona \\
Sul piano politico e sociale funziona meno \t si era tentata una costituzione europea, e su quali fondamenti comuni si legassero tutti i paesi \t ma fallisce \\
Stati europei non si riconoscono in documento comune \t gli stati sono infatti reticenti di trasferire alcuni loro poteri a strutture sovranazionali
\end{document}
